\section{Simulation design}
\label{sec:design}

Experiment~4 is a paired recovery comparison.  For each task, the
data-generating POMP simulates an epidemic and reported cases once.  The
resulting observation times and case counts are saved and supplied unchanged
to the Gamma-noise and constant-\(B\) fitting POMPs.  Pairing removes
between-data-set variation from the model contrast: both fits must explain the
same realized epidemic reports.

Within a task, the deterministic transmission path drives a new stochastic
SIR realization, the accumulator records infections between daily
observations, and the negative-binomial model produces one case count at each
observation time.  Only these observation times and counts enter the two
fitting workflows.  The unobserved generating states are retained for
evaluation but are not supplied to IF2 or to the final filters.

The data-generating transmission path is
\begin{equation}
\Btrue(t)
=
\begin{cases}
4, & 0\leq t<5,\\[3pt]
2, & 5\leq t\leq10,
\end{cases}
\qquad \text{\(t\) in weeks.}
\label{eq:truth}
\end{equation}
This path is deterministic.  Neither fitting model receives
\(\Btrue(t)\), the values in Equation~\ref{eq:truth}, or the week-5 change
point.  The principal settings are summarized in
Table~\ref{tab:settings}.

The 70 daily observations span the ten-week horizon; the scheduled
observation at week~5 is evaluated under the post-change rate.  The Euler grid
is finer than the observation grid, with approximately four
process updates between successive daily measurements.  Fixing the recovery,
reporting, population, and initial-state quantities at their generating values
makes this a focused transmission-recovery experiment: it avoids adding
nuisance-parameter confounding to the primary comparison.

\begin{table}[htbp]
\centering
\caption{Experiment~4 simulation settings.}
\label{tab:settings}
\begin{tabularx}{0.88\textwidth}{>{\raggedright\arraybackslash}p{0.43\textwidth}X}
\toprule
Setting & Value \\
\midrule
Paired accepted data sets & 200 \\
Acceptance criterion & \(\max(H)>20\) \\
Duration & 10 weeks \\
Observation schedule & 70 daily observations, \(1/7,\ldots,10\) weeks \\
Euler step & \(\Delta t=1/30\) week \\
Prescribed transmission path & \(4\) before week~5; \(2\) from week~5 onward \\
Population size & \(N=10000\) \\
Recovery rate & \(\mu_{IR}=3~\mathrm{week}^{-1}\) \\
Reporting fraction & \(\rho=0.5\) \\
Negative-binomial size & \(k=10\) \\
Initial state & \(S(0)=9990,\ I(0)=10,\ R(0)=H(0)=0\) \\
\bottomrule
\end{tabularx}
\end{table}

The acceptance rule conditions the comparison on simulated outbreaks with
more than 20 accumulated infections in at least one observation interval.
Consequently, the reported recovery results describe accepted outbreaks under
these settings rather than all trajectories that the data-generating POMP
could produce.  Experiments~1--3 were developmental steps used to build and
calibrate the workflow; only Experiment~4 supplies the main scientific
evidence.  Their roles are summarized in
Appendix~\ref{app:developmental}.
